% Shared pleading-paper preamble for CGC-25-631802 TAC-RESPONSE court PDFs.
\documentclass{article}
\usepackage[utf8]{inputenc}
\usepackage{graphicx}
\usepackage{geometry}
    \geometry{top=1in, bottom=2.2in, left=1.0in, right=1.0in, textheight=336pt}
    \setlength{\footskip}{55pt}
\usepackage{ulem}
\usepackage{fancyhdr}
    \renewcommand{\headrulewidth}{0pt}
    \renewcommand{\footrulewidth}{0pt}
\usepackage{parskip}
    \setlength{\parskip}{0mm}
\usepackage{quoting}
    \quotingsetup{leftmargin=0.5in, rightmargin=0.5in, vskip=-1.5mm}
    \makeatletter
        \g@addto@macro\quoting\singlespacing
        \g@addto@macro\quoting{\vspace{-2mm}}
    \makeatother
\usepackage{mathptmx}
\usepackage{amssymb}
\renewcommand{\rmdefault}{ptm}
\renewcommand{\normalsize}{\fontsize{12}{12}\selectfont}
\newcommand*{\ptm}{\fontfamily{ptm} \selectfont \fontsize{12}{12}\selectfont}
\usepackage{setspace}
    \doublespacing
\newcommand{\tab}{\hspace*{.0in}}
\newenvironment{tightcenter}{%
    \setlength\topsep{0pt}
    \setlength\parskip{0pt}
    \begin{center}
}{%
    \end{center}
}
\usepackage{tikz}
\usetikzlibrary{calc}
\AddToHook{shipout/background}{%
    \begin{tikzpicture}[remember picture, overlay]
        \draw[line width=0.5pt] ($(current page.north west) + (0.75in,0)$) -- ($(current page.south west) + (0.75in,0)$);
        \draw[line width=0.5pt] ($(current page.north west) + (0.80in,0)$) -- ($(current page.south west) + (0.80in,0)$);
    \end{tikzpicture}%
}
\usepackage[left, pagewise]{lineno}
\setlength{\linenumbersep}{0.35in}
\renewcommand{\linenumberfont}{\normalfont\small}
\usepackage{titlesec}
\usepackage[hidelinks]{hyperref}
\usepackage{bookmark}
\usepackage{textcomp}
\usepackage{longtable}
\usepackage{booktabs}
\usepackage{array}
\usepackage{multirow}
\usepackage{calc}
% Pandoc pipe-table column widths use \real{#}; map to pgf math (pgf loaded via tikz).
\newcommand{\real}[1]{\pgfmathparse{#1}\pgfmathresult}
\titleformat{\section}{\fontsize{15}{15}\selectfont\normalfont\bfseries\uppercase}{}{0em}{}
\titleformat{\subsection}{\fontsize{15}{15}\selectfont\normalfont\bfseries\uppercase}{}{0em}{}
\titleformat{\subsubsection}{\normalfont\bfseries}{}{0em}{}
\titlespacing*{\section}{0pt}{0pt}{0pt}
\titlespacing*{\subsection}{0pt}{0pt}{0pt}
\titlespacing*{\subsubsection}{0pt}{0pt}{0pt}
\newcommand{\calrules}[1]{Cal. Rules of Court, rule #1}
\newcommand{\ccp}[1]{Code Civ. Proc., \textsection~#1}
\newcommand{\evidcode}[1]{Evid. Code, \textsection~#1}
\newcommand{\vehcode}[1]{Veh. Code, \textsection~#1}
\providecommand{\tightlist}{%
  \setlength{\itemsep}{0pt}\setlength{\parskip}{0pt}}
% --- TOC / TOA hyperlink infrastructure ---
\newcommand{\tacctoclink}[2]{\hyperref[toc:#1]{#2}}
% In-text citation link to authority label (target is usually TABLE OF AUTHORITIES row).
\newcommand{\citecase}[2]{\textit{\hyperref[auth:#1]{#2}}}
\newcommand{\citestatute}[2]{\hyperref[auth:#1]{#2}}
\newcommand{\toclink}[2]{\hyperref[toc:#1]{#2}}
% Authority entries in TABLE OF AUTHORITIES (#1 label slug, #2 italic case name or title, #3 citation tail)
\newcommand{\caseentry}[3]{\par\noindent\hangindent=2em\hangafter=1%
  \phantomsection\label{auth:#1}\textit{#2}, #3\par}
\newcommand{\statuteentry}[3]{\par\noindent\hangindent=2em\hangafter=1%
  \phantomsection\label{auth:#1}#2\ #3\par}
\newcommand{\ruleentry}[2]{\par\noindent\hangindent=2em\hangafter=1%
  \phantomsection\label{auth:#1}#2\par}
\makeatletter
\newlength{\CourtTOCmarkwd}
\newlength{\CourtTOCpagewd}
\newlength{\CourtTOCtitlewd}
\setlength{\CourtTOCmarkwd}{2.8em}
\setlength{\CourtTOCpagewd}{4.2em}
\newcommand{\CourtTOCRow}[3]{%
  \begingroup
  \setlength{\CourtTOCtitlewd}{\dimexpr\linewidth - \CourtTOCmarkwd - \CourtTOCpagewd - 0.6em\relax}%
  \par\nobreak\noindent
  \begin{minipage}[b]{\CourtTOCmarkwd}\raggedright\strut #1\strut\end{minipage}\hspace{0.3em}%
  \begin{minipage}[b]{\CourtTOCtitlewd}\raggedright\strut #2\strut\end{minipage}\hspace{0.3em}%
  \begin{minipage}[b]{\CourtTOCpagewd}\raggedleft\strut #3\strut\end{minipage}\par
  \endgroup
  \vspace{1pt}%
}
\newcommand{\CourtTOCSubitem}[2]{\CourtTOCRow{}{#1}{#2}}
\newcommand{\CourtTOCSubsubitem}[2]{\CourtTOCRow{}{\quad\quad #1}{#2}}
\makeatother

\begin{document}
\nolinenumbers
\begin{singlespace}
Franciscus Dylan Rosario, \\
7624 Melody \\
Rohnert Park, CA 94928 \\
E: soltrinox@gmail.com \\
Plaintiff, In Pro Per \\
\end{singlespace}
\vspace*{9mm}
\begin{tightcenter}
\textbf{SUPERIOR COURT OF THE STATE OF CALIFORNIA} \\
\textbf{COUNTY OF SAN FRANCISCO}
\end{tightcenter}
\vspace*{2.25mm}
\nolinenumbers
\begin{minipage}[t]{3in}
    \raggedright
    \textbf{Franciscus Dylan Rosario,} \\ \textbf{PLAINTIFF}, \\
    \hspace{1cm} v. \\
    \small
    CSAA Insurance Exchange; \\
    and DOES 1 through 50, inclusive; \\
    \normalsize
    \textbf{DEFENDANTS}
    \makebox[3in]{\hrulefill}
\end{minipage}%
\hspace{3mm}
\begin{minipage}[t]{0.5pt}
    \vspace{0pt}
    \rule{0.5pt}{2.0in}
\end{minipage}
\hspace{3mm}
\begin{minipage}[t]{3in}
    \raggedright
    \noindent \textbf{Case No.: CGC-25-631802} \\
    ~\\
    \textbf{CURE MATRIX --- DEFENSE OBJECTION MAP}
\end{minipage}
\vspace*{6mm}
\linenumbers
\section{Second-Layer Cure Matrix If SAC Defects Survive}\label{second-layer-cure-matrix-if-sac-defects-survive}

\textbf{Case No.:} CGC-25-631802\\
\textbf{Subject:} Third Amended Complaint for Fraud and UCL\\
\textbf{Status:} CONTINGENT --- filed only if Round D defects surface after SAC ruling\\
\textbf{Purpose:} Maps each defense objection from the April 22, 2026 CCP \textsection{} 436 Motion to Strike to the specific TAC paragraph(s) that cure it.

\textbf{First Cure:} SAC-leave motion (hearing Apr 30, 2026)\\
\textbf{Second Cure:} This TAC packet (hearing TBD on OST if needed)

\textbf{Underlying Action:} Rosario v. Abdelhalim, Case No.~CGC-21-594102

\begin{center}\rule{0.5\linewidth}{0.5pt}\end{center}

\section{Master Cure Matrix --- Seven Defense Objections}\label{master-cure-matrix-seven-defense-objections}

{\def\LTcaptype{table} % do not increment counter
\begin{longtable}[]{@{}
  >{\raggedright\arraybackslash}p{(\linewidth - 12\tabcolsep) * \real{0.0250}}
  >{\raggedright\arraybackslash}p{(\linewidth - 12\tabcolsep) * \real{0.1583}}
  >{\raggedright\arraybackslash}p{(\linewidth - 12\tabcolsep) * \real{0.1583}}
  >{\raggedright\arraybackslash}p{(\linewidth - 12\tabcolsep) * \real{0.0833}}
  >{\raggedright\arraybackslash}p{(\linewidth - 12\tabcolsep) * \real{0.1333}}
  >{\raggedright\arraybackslash}p{(\linewidth - 12\tabcolsep) * \real{0.1250}}
  >{\raggedright\arraybackslash}p{(\linewidth - 12\tabcolsep) * \real{0.3167}}@{}}
\toprule\noalign{}
\begin{minipage}[b]{\linewidth}\raggedright
\#
\end{minipage} & \begin{minipage}[b]{\linewidth}\raggedright
Defense Objection
\end{minipage} & \begin{minipage}[b]{\linewidth}\raggedright
Defense MPA Basis
\end{minipage} & \begin{minipage}[b]{\linewidth}\raggedright
TAC Cure
\end{minipage} & \begin{minipage}[b]{\linewidth}\raggedright
TAC Section/¶¶
\end{minipage} & \begin{minipage}[b]{\linewidth}\raggedright
Key Authority
\end{minipage} & \begin{minipage}[b]{\linewidth}\raggedright
Cal. Supreme / Additional Authority
\end{minipage} \\
\midrule\noalign{}
\endhead
\bottomrule\noalign{}
\endlastfoot
1 & \emph{Trope} --- pro-per fees barred & Pro se cannot recover fees & Deletion of pro-per fees; \emph{Prentice} damages element with four line-item categories (incl.~\$75k cost-judgment as damages) & \textsection{}III, ¶¶ 22--24, 24A & \emph{Prentice v. North American Title} (1963) 59 Cal.2d 618, 620 & \emph{Prentice} 59 Cal.2d 618 (Cal. Sup. Ct.) \\
2 & Oppression / malice not pleaded & Civ. Code \textsection{} 3294 & \emph{Lazar} particulars; deny-delay-defend pattern; confession/victim/eyewitness triple; pleading-vs-evidentiary bright line & \textsection{}IV, ¶¶ 25--30, 25A & \emph{Lazar} (1996) 12 Cal.4th 631; \emph{American Airlines} (2018) 28 Cal.App.5th 1077 & \emph{G.D. Searle \& Co.~v. Superior Court} (1975) 49 Cal.App.3d 22, 29 \\
3 & Fraud lacks particularity & \emph{Lazar} requirements & Sarah Rash named; Claim No.~1004-09-1054; Feb.~25, 2021; specific representations; evidence triple with SFPD incident no. and first-responder identifiers & \textsection{}I, ¶¶ 3--4, 8; \textsection{}IV, ¶ 25; ¶ 29 & \emph{Lazar} & \emph{Lazar} 12 Cal.4th 631 (Cal. Sup. Ct.) \\
4 & Litigation privilege / \emph{Ribas} & Civ. Code \textsection{} 47(b) & Pre-litigation carveout (three-factor Action Apartment walk-through); Flatley illegality carve-out; non-communicative conduct; \emph{Ribas} distinction (fraud $\neq$ IIED) & \textsection{}II, ¶¶ 14--21, 16A, 21A & \emph{Action Apartment} (2007) 41 Cal.4th 1232; \emph{Rusheen} (2006) 37 Cal.4th 1048 & \emph{Flatley v. Mauro} (2006) 39 Cal.4th 299, 317--21 \\
5 & Standing / \emph{Shaolian} & No direct duty to claimant & Direct-addressee theory; voluntary direct communication; Civ. Code \textsection{} 2338 agency / ratification spine closing the ``CSAA did not itself act'' imputation gap & \textsection{}I.5, ¶¶ 5--6, 6A & --- & Civ. Code \textsection{} 2338; \emph{Doctors' Company v. Superior Court} (1989) 49 Cal.3d 39 \\
6 & Collateral estoppel / \textsection{} 1032 & PI damages already adjudicated & Distinct primary right; express PI disclaimer; 2021-forward anchor; Lucido five-factor walk-through; 801/802 double-recovery reservation & \textsection{}VI, ¶¶ 41--44, 41A, 44A & --- & \emph{Lucido v. Superior Court} (1990) 51 Cal.3d 335; \emph{Mycogen Corp.~v. Monsanto Co.} (2002) 28 Cal.4th 888; \emph{Crowley v. Katleman} (1994) 8 Cal.4th 666 \\
7 & Investigation standard not pleaded & Title 10 CCR & Title 10 CCR \textsection{} 2695.7 requirements; itemized failures; \emph{Moradi-Shalal} carve-out; Ins. Code \textsection{} 790.03(h)(5) liability-reasonably-clear UCL predicate & \textsection{}VII, ¶¶ 45--47, 47A & 10 CCR \textsection{} 2695.4, \textsection{} 2695.7; \emph{State Farm} (1996) 45 Cal.App.4th 1093 & Ins. Code \textsection{} 790.03(h)(5); \emph{Moradi-Shalal} (1988) 46 Cal.3d 287 \\
8 & Anti-SLAPP / \textsection{} 425.16 overlap with \textsection{} 436 & Litigation conduct claim & Express wrongful-conduct-scope disavowal of all litigation-conduct predicates; anchors on 2021 pre-litigation claim handling only & First CoA, ¶ 52A & CCP \textsection{} 425.16(e); \emph{Flatley v. Mauro} & \emph{Flatley v. Mauro} 39 Cal.4th 299 \\
9 & ``Ultimate admission cures fraud'' / UCL fraudulent-prong sufficiency & No reasonable-consumer deception & Reasonable-consumer test pleaded as UCL fraudulent-prong standard & Third CoA, ¶¶ 66, 66A & Bus. \& Prof.~Code \textsection{} 17200 & \emph{Williams v. Gerber Products Co.} (9th Cir. 2008) 552 F.3d 934, 938; \emph{Chapman v. Skype, Inc.} (2013) 220 Cal.App.4th 217, 226 \\
\end{longtable}
}

\begin{center}\rule{0.5\linewidth}{0.5pt}\end{center}

\section{Detailed Cure Analysis}\label{detailed-cure-analysis}

\subsection{\texorpdfstring{Objection 1: \emph{Trope} --- Pro-Per Fees Barred}{Objection 1: Trope --- Pro-Per Fees Barred}}\label{objection-1-trope-pro-per-fees-barred}

\textbf{Defense Argument:} Pro se litigant cannot recover attorney's fees under \emph{Trope v. Katz}.

\textbf{TAC Cure:}

\begin{itemize}
\tightlist
\item
  \textbf{¶ 22:} Express deletion of pro-per fee demands
\item
  \textbf{¶¶ 23--24:} \emph{Prentice} fees pleaded as tort damages element (three-element test):

  \begin{itemize}
  \tightlist
  \item
    Element 1: Fraud
  \item
    Element 2: Induced litigation
  \item
    Element 3: Fees recoverable as damages, not fee award
  \end{itemize}
\item
  \textbf{¶ 24A:} Four-line-item \emph{Prentice} enumeration --- (1) professional fees/costs incurred in CGC-21-594102; (2) the \$75,000 cost judgment as induced-litigation damages (re-framing MPA2 footnote 2 rhetoric into Plaintiff's damages element); (3) transcript/record/appeal preparation costs; (4) investigative/records-acquisition costs (911, CAD, BWC) that CSAA itself should have obtained.
\end{itemize}

\textbf{Authority:} \emph{Trope v. Katz} (1995) 11 Cal.4th 274; \emph{Prentice v. North American Title Guaranty Corp.} (1963) 59 Cal.2d 618, 620

\begin{center}\rule{0.5\linewidth}{0.5pt}\end{center}

\subsection{Objection 2: Oppression / Malice Not Pleaded}\label{objection-2-oppression-malice-not-pleaded}

\textbf{Defense Argument:} Punitive damages allegations fail to meet \textsection{} 3294 standard.

\textbf{TAC Cure:}

\begin{itemize}
\tightlist
\item
  \textbf{¶¶ 25--27:} \emph{Lazar} particulars block with \textsection{} 3294(b) managing-agent ratification
\item
  \textbf{¶ 25A:} Pleading-stage vs.~trial-stage bright line --- \emph{G.D. Searle \& Co.~v. Superior Court} (1975) 49 Cal.App.3d 22, 29: at pleading, plaintiff need only allege facts which, \textbf{if proven by clear and convincing evidence at trial}, would authorize punitive recovery. Forecloses the anticipated defense misuse of \emph{In re Angelia P.} (1981) 28 Cal.3d 908 as a pleading standard.
\item
  \textbf{¶¶ 28--30:} Deny-delay-defend institutional pattern; ``despicable conduct''
\item
  \textbf{¶ 29:} Confession/victim report/eyewitness corroboration triple --- claim denied despite all three (refined with SFPD Incident Report No.~210083458 (conformed to proof), named eyewitness Justin Rosemond, and public-records/subpoena availability within days of February 4, 2021).
\end{itemize}

\textbf{Authority:} Civ. Code \textsection{} 3294; \emph{Lazar}; \emph{American Airlines, Inc.~v. Sheppard, Mullin} (2018) 28 Cal.App.5th 1077; \emph{G.D. Searle \& Co.~v. Superior Court} (1975) 49 Cal.App.3d 22, 29

\begin{center}\rule{0.5\linewidth}{0.5pt}\end{center}

\subsection{Objection 3: Fraud Lacks Particularity}\label{objection-3-fraud-lacks-particularity}

\textbf{Defense Argument:} Complaint fails \emph{Lazar} who/what/when/where/how requirements.

\textbf{TAC Cure:}

\begin{itemize}
\tightlist
\item
  \textbf{WHO:} Sarah Rash, Claims Representative (\textsection{}I, ¶ 3; \textsection{}IV, ¶ 25(a))
\item
  \textbf{WHAT:} False representation of investigative completeness (\textsection{}I, ¶ 4)
\item
  \textbf{WHEN:} February 25, 2021 --- 21 days after collision (\textsection{}I, ¶ 3; \textsection{}IV, ¶ 25(c))
\item
  \textbf{WHERE:} CSAA claims office; communication to Plaintiff in San Francisco (\textsection{}IV, ¶ 25(d))
\item
  \textbf{HOW:} Written denial letter without investigation (\textsection{}IV, ¶ 25(e))
\item
  \textbf{BY WHAT MEANS:} U.S. Mail (\textsection{}IV, ¶ 25(f))
\item
  \textbf{Claim No.:} 1004-09-1054 (\textsection{}I, ¶ 3)
\end{itemize}

\textbf{Authority:} \emph{Lazar v. Superior Court} (1996) 12 Cal.4th 631

\begin{center}\rule{0.5\linewidth}{0.5pt}\end{center}

\subsection{\texorpdfstring{Objection 4: Litigation Privilege / \emph{Ribas}}{Objection 4: Litigation Privilege / Ribas}}\label{objection-4-litigation-privilege-ribas}

\textbf{Defense Argument:} Claims barred by Civil Code \textsection{} 47(b) and \emph{Ribas v. Clark}.

\textbf{TAC Cure:}

\begin{itemize}
\tightlist
\item
  \textbf{¶¶ 14--16:} Pre-litigation carveout --- denial issued before any litigation
\item
  \textbf{¶ 16A:} Three-factor \emph{Action Apartment} walk-through --- affirmative pleading that, as of February 25, 2021, \textbf{none} of the three contemplation factors was present in CSAA's favor: (a) no litigation pending or contemplated; (b) no counsel retained, no demand exchanged, no tolling agreement; (c) no litigation-hold / preservation letter and no litigation reserve.
\item
  \textbf{¶¶ 17--19:} \emph{Ribas} distinction --- \emph{Ribas} addresses IIED, not fraud
\item
  \textbf{¶¶ 20--21:} Non-communicative conduct (failure to investigate) outside privilege
\item
  \textbf{¶ 21A:} \emph{Flatley v. Mauro} (2006) 39 Cal.4th 299, 317--21, illegality carve-out --- categorical removal of \textsection{} 47(b) for conduct that is illegal as a matter of law, including fraud under Civ. Code \textsection{}\textsection{} 1709, 1710.
\end{itemize}

\textbf{Authority:} - \emph{Action Apartment Ass'n v. City of Santa Monica} (2007) 41 Cal.4th 1232, 1248--49 - \emph{Flatley v. Mauro} (2006) 39 Cal.4th 299, 317--21 - \emph{Rusheen v. Cohen} (2006) 37 Cal.4th 1048 - \emph{Ribas v. Clark} (1985) 38 Cal.3d 355 (distinguished)

\begin{center}\rule{0.5\linewidth}{0.5pt}\end{center}

\subsection{\texorpdfstring{Objection 5: Standing / \emph{Shaolian}}{Objection 5: Standing / Shaolian}}\label{objection-5-standing-shaolian}

\textbf{Defense Argument:} No direct duty to third-party claimant under \emph{Shaolian v. Safeco}.

\textbf{TAC Cure:}

\begin{itemize}
\tightlist
\item
  \textbf{¶ 5:} Direct-addressee theory --- Plaintiff is the person to whom the representation was made
\item
  \textbf{¶ 6:} CSAA voluntarily undertook direct communicative relationship with unrepresented claimant
\item
  \textbf{¶ 6A:} Civil Code \textsection{} 2338 agency/ratification spine --- CSAA is directly liable for wrongful acts and willful omissions of Sarah Rash, claims department personnel, and retained claims professionals acting within scope; ratification pleaded through failure to correct after the January 17, 2025 / February 18, 2025 / April 16, 2025 record events. Anchored in \emph{Doctors' Company v. Superior Court} (1989) 49 Cal.3d 39 (insurer directly liable when it acts outside ordinary insurer role).
\end{itemize}

\textbf{Distinction:} \emph{Shaolian} addresses policy benefits claims; this is a common law fraud claim by the direct addressee of the misrepresentation, with CSAA directly liable under Civ. Code \textsection{} 2338 for the agent-acts by which the misrepresentation was made.

\begin{center}\rule{0.5\linewidth}{0.5pt}\end{center}

\subsection{Objection 6: Collateral Estoppel / \textsection{} 1032}\label{objection-6-collateral-estoppel-1032}

\textbf{Defense Argument:} Damages duplicate those sought or awarded in CGC-21-594102.

\textbf{TAC Cure:}

\begin{itemize}
\tightlist
\item
  \textbf{¶ 41:} Distinct primary right --- freedom from misrepresentation about investigative completeness
\item
  \textbf{¶ 41A:} \emph{Lucido} five-factor walk-through against CGC-21-594102 record: (i) identical issue --- no (negligence vs.~fraud); (ii) actually litigated --- no (CSAA not a party; no fraud tendered); (iii) necessarily decided --- no; (iv) final judgment --- no (pending on appeal A173827); (v) same party or privy --- no. Primary-rights framework under \emph{Mycogen Corp.~v. Monsanto Co.} (2002) 28 Cal.4th 888, 904, and \emph{Crowley v. Katleman} (1994) 8 Cal.4th 666.
\item
  \textbf{¶ 43:} Express disclaimer of PI damages
\item
  \textbf{¶ 44:} 2021-forward damages anchor --- all damages arise from Feb.~25, 2021 denial onward
\item
  \textbf{¶ 44A:} Express 801/802 double-recovery reservation --- no overlap between legal-damages relief here and equitable-vacatur relief in CGC-25-631801; pre-judgment election of remedies preserved.
\end{itemize}

\textbf{Authority:} \emph{Lucido v. Superior Court} (1990) 51 Cal.3d 335, 341; \emph{Mycogen Corp.~v. Monsanto Co.} (2002) 28 Cal.4th 888, 904; \emph{Crowley v. Katleman} (1994) 8 Cal.4th 666, 681--82.

\begin{center}\rule{0.5\linewidth}{0.5pt}\end{center}

\subsection{Objection 7: Investigation Standard Not Pleaded}\label{objection-7-investigation-standard-not-pleaded}

\textbf{Defense Argument:} Complaint fails to allege what investigation was required.

\textbf{TAC Cure:}

\begin{itemize}
\tightlist
\item
  \textbf{¶ 45:} What thorough investigation requires under 10 CCR \textsection{} 2695.7(b)
\item
  \textbf{¶ 46:} Itemized failures --- CSAA did not obtain 911, BWC, CAD records
\item
  \textbf{¶ 47:} \emph{Moradi-Shalal} carve-out --- not a private \textsection{} 790.03 claim; UCL predicate only
\item
  \textbf{¶ 47A:} Ins. Code \textsection{} 790.03(h)(5) ``liability reasonably clear'' pleaded \textbf{solely} as a UCL unlawful-prong predicate consistent with the \emph{Moradi-Shalal} carve-out; liability was reasonably clear to CSAA on the face of SFPD Incident Report No.~210083458 (conformed to proof) and through routine public-records availability of 911/CAD/BWC within days of the February 4, 2021 collision.
\end{itemize}

\textbf{Authority:} 10 CCR \textsection{}\textsection{} 2695.4, 2695.7(b); Ins. Code \textsection{} 790.03(h)(5); \emph{Moradi-Shalal v. Fireman's Fund Ins. Cos.} (1988) 46 Cal.3d 287; \emph{State Farm Fire \& Casualty Co.~v. Superior Court} (1996) 45 Cal.App.4th 1093, 1103--04.

\begin{center}\rule{0.5\linewidth}{0.5pt}\end{center}

\subsection{Objection 8: Anti-SLAPP Overlap / CCP \textsection{} 425.16}\label{objection-8-anti-slapp-overlap-ccp-425.16}

\textbf{Defense Argument (anticipated):} Claims are predicated on protected litigation conduct within \textsection{} 425.16(e).

\textbf{TAC Cure:}

\begin{itemize}
\tightlist
\item
  \textbf{¶ 52A:} Express wrongful-conduct-scope disavowal --- no liability predicated on filing/prosecution of CGC-21-594102, in-limine practice, or statements made in connection with a judicial proceeding. All wrongful conduct anchored in 2021 pre-litigation claim handling, already placed outside \textsection{} 47(b) by ¶¶ 16A and 21A. Collapses \textsection{} 436 Cluster B overlap with \textsection{} 425.16 step-1.
\end{itemize}

\textbf{Authority:} CCP \textsection{} 425.16(e); \emph{Flatley v. Mauro} (2006) 39 Cal.4th 299, 317--21.

\begin{center}\rule{0.5\linewidth}{0.5pt}\end{center}

\subsection{Objection 9: Reasonable-Consumer / UCL Fraudulent-Prong Sufficiency}\label{objection-9-reasonable-consumer-ucl-fraudulent-prong-sufficiency}

\textbf{Defense Argument (anticipated):} UCL fraudulent-prong claim fails absent allegations of likely deception of a reasonable consumer.

\textbf{TAC Cure:}

\begin{itemize}
\tightlist
\item
  \textbf{¶¶ 66, 66A:} Reasonable-consumer test pleaded --- a business practice is ``fraudulent'' under \textsection{} 17200 if members of the public are likely to be deceived. Representation of ``concluded investigation'' to unrepresented third-party claimants who lack means to audit an insurer's claim file is plausibly likely to deceive a reasonable claimant.
\end{itemize}

\textbf{Authority:} Bus. \& Prof.~Code \textsection{} 17200; \emph{Williams v. Gerber Products Co.} (9th Cir. 2008) 552 F.3d 934, 938; \emph{Chapman v. Skype, Inc.} (2013) 220 Cal.App.4th 217, 226.

\begin{center}\rule{0.5\linewidth}{0.5pt}\end{center}

\section{Cross-Reference: TAC Sections to Defense Objections}\label{cross-reference-tac-sections-to-defense-objections}

{\def\LTcaptype{table} % do not increment counter
\begin{longtable}[]{@{}
  >{\raggedright\arraybackslash}p{(\linewidth - 4\tabcolsep) * \real{0.4048}}
  >{\raggedright\arraybackslash}p{(\linewidth - 4\tabcolsep) * \real{0.1667}}
  >{\raggedright\arraybackslash}p{(\linewidth - 4\tabcolsep) * \real{0.4286}}@{}}
\toprule\noalign{}
\begin{minipage}[b]{\linewidth}\raggedright
TAC Section / ¶
\end{minipage} & \begin{minipage}[b]{\linewidth}\raggedright
Title
\end{minipage} & \begin{minipage}[b]{\linewidth}\raggedright
Cures Objections
\end{minipage} \\
\midrule\noalign{}
\endhead
\bottomrule\noalign{}
\endlastfoot
I & Direct-Duty Gravamen & 3, 5 \\
I.5 & Direct-Addressee Theory & 5 \\
¶ 6A & Civ. Code \textsection{} 2338 Agency / Ratification & 5 \\
¶ 8 & Evidence triple with SFPD identifiers & 3 \\
I.6 & Gravamen Disclaimer & --- \\
II & Pre-Litigation Privilege Carveout & 4 \\
¶ 16A & Action Apartment three-factor absence & 4 \\
¶ 21A & Flatley illegality carve-out & 4 \\
III & Attorney's Fees Cure & 1 \\
¶ 24A & Prentice four-line-item enumeration (incl.~\$75k cost-judgment as damages) & 1 \\
IV & Punitive Damages Specificity & 2, 3 \\
¶ 25A & Pleading-vs-evidentiary bright line (G.D. Searle) & 2 \\
¶ 29 & Evidence triple refined with SFPD identifiers & 2, 3 \\
V & Damages & --- \\
VI & Primary Rights / Reservation & 6 \\
¶ 41A & Lucido five-factor walk-through & 6 \\
¶ 44A & 801/802 double-recovery reservation & 6 \\
VII & Title 10 Investigation Standard & 7 \\
¶ 47A & Ins. Code \textsection{} 790.03(h)(5) UCL predicate & 7 \\
VIII & Delayed Discovery & --- \\
First CoA ¶ 52A & Anti-SLAPP wrongful-conduct-scope disavowal & 8 \\
Third CoA ¶ 66A & Reasonable-consumer test & 9 \\
\end{longtable}
}

\begin{center}\rule{0.5\linewidth}{0.5pt}\end{center}

\section{Uncured Gaps --- Closed by TAC}\label{uncured-gaps-closed-by-tac}

The prior ``CSAA did not itself act'' / trial-conduct imputation gap is \textbf{closed} by new ¶ 6A (Civ. Code \textsection{} 2338 agency and ratification spine). CSAA is directly and expressly liable for the wrongful acts and willful omissions of Sarah Rash, claims department personnel, and retained claims professionals within the scope of the claims-handling agency; ratification is separately pleaded through CSAA's failure to correct the February 25, 2021 denial after the January 17, 2025 / February 18, 2025 / April 16, 2025 record events described in TAC Section VIII. See \emph{Doctors' Company v. Superior Court} (1989) 49 Cal.3d 39. No open uncured gaps remain.

\begin{center}\rule{0.5\linewidth}{0.5pt}\end{center}

\section{Validation Checklist}\label{validation-checklist}

\begin{itemize}
\tightlist
\item[$\boxtimes$]
  All seven defense objections mapped to TAC cure
\item[$\boxtimes$]
  Each cure supported by legal authority and/or factual specificity
\item[$\boxtimes$]
  Sarah Rash and Claim No.~1004-09-1054 filled throughout
\item[$\boxtimes$]
  Three-anchor delayed-discovery block included (Jan.~17 / Feb.~18 / Apr.~16, 2025)
\item[$\boxtimes$]
  Deny-delay-defend institutional framing included
\item[$\boxtimes$]
  Direct-addressee theory as \emph{Shaolian} rebuttal included
\item[$\boxtimes$]
  UCL injunction tightened to ``issuing denial letters stating investigation is concluded''
\item[$\boxtimes$]
  SAC ¶6A primary-right statement included
\item[$\boxtimes$]
  Causes of action restructured to three-count architecture (Fraud-Omission/Extrinsic; Concealment/Half-Truth; UCL \textsection{} 17200)
\item[$\boxtimes$]
  \emph{Moradi-Shalal} carve-out included
\item[$\boxtimes$]
  Pro-per fee demands deleted
\end{itemize}

\begin{center}\rule{0.5\linewidth}{0.5pt}\end{center}

\emph{Last updated: April 22, 2026}

\nolinenumbers
% Generated by build_tac_response_802_court_pdfs.write_tac_signature_block_tex()
\vspace{0.5\baselineskip}
\begin{singlespace}
\begin{flushleft}
Dated: April 22, 2026 \\
\vspace{0.25in}
\noindent\includegraphics[height=0.5in]{../../../Support/signature.png}\\[\baselineskip]
FRANCISCUS DYLAN ROSARIO \\
Plaintiff, In Pro Per \\
\vspace{0.15in}
7624 Melody \\
Rohnert Park, CA 94928 \\
650.488.7510 \\
E: soltrinox@gmail.com
\end{flushleft}
\end{singlespace}


\end{document}
